\bookchapter{Capítulo 35: Niños y móviles — Guía educativa para la era digital}

\chapterintro{: El objetivo se ensucia con el uso y las fotos salen borrosas. Límpialo con un paño suave antes de tomar fotos importantes.}

\booksection{Usa la cuadrícula para la composición}

: Ajustes > Cámara > Cuadrícula. La base de la composición con la regla de tercios; te ayuda a encontrar rápidamente el centro visual.

\begin{itemize}
  \item \textbf{Bloquea el enfoque y la exposición}: Mantén pulsado el sujeto en el encuadre durante 2 segundos; aparecerá "Bloqueo AE/AF". Así, al grabar o tomar fotos, la imagen no cambiará de brillo al moverte.\textbf{No uses el zoom digital}: El zoom digital solo recorta la imagen y pierde calidad. Acércate físicamente al sujeto o usa el zoom óptico si el teléfono lo tiene.
  \item \textbf{Usa el HDR para escenas con mucho contraste}: El HDR (High Dynamic Range) mejora las fotos con mucho contraste (cielo brillante y sombras oscuras). Actívalo en Ajustes > Cámara > HDR automático.
  \item \textbf{No uses el flash a menos que sea necesario}: El flash del teléfono es potente pero produce una luz dura y poco natural. Úsalo solo en la oscuridad total; en condiciones de poca luz, apoya el teléfono o usa el modo nocturno.
  \item \textbf{Usa el modo retrato para personas}: El modo retrato desenfoca el fondo y hace que la persona destaque. Úsalo para retratos de personas; no lo uses para paisajes ni objetos.
\end{itemize}

\begin{tcbinfo}{Usa el modo nocturno para fotos de noche}
: El modo nocturno del teléfono toma varias fotos y las combina para mejorar la luminosidad. Úsalo en escenas de poca luz; mantén el teléfono estable.

\end{tcbinfo}

\booksubsection{Edita las fotos antes de compartirlas}

: Usa la app de Fotos para editar: recorta, ajusta el brillo y el contraste, aplica filtros. Una edición rápida mejora mucho la foto.

\begin{itemize}
  \item \textbf{Haz copia de seguridad de tus fotos}: Las fotos son recuerdos irreemplazables. Activa la copia de seguridad automática (iCloud / Google Photos / copia local del fabricante) para no perderlas.
  \item \textbf{X. Diez consejos para el almacenamiento}El almacenamiento del teléfono se llena rápidamente con fotos, vídeos y apps. Estos diez consejos te ayudan a gestionarlo:
  \item \textbf{Desinstala las apps que no uses}: Cada app instalada consume almacenamiento. Desinstala las apps que no uses; puedes reinstalarlas cuando las necesites.
\end{itemize}

Borra las fotos y vídeos duplicados\textbf{: Usa la función “Duplicados” de la app de Fotos (iPhone) o apps de terceros para encontrar y borrar fotos y vídeos duplicados.}Vacía la caché de las apps

\booksection{: Ajustes > Apps > [app] > Almacenamiento > Borrar caché. Esto libera espacio sin borrar datos de la app.}

Borra los mensajes antiguos

\booksubsection{: Los mensajes de WeChat y SMS ocupan espacio. Borra los mensajes antiguos que no necesites; en WeChat, usa “Gestión de almacenamiento” para limpiar.}

\begin{itemize}
  \item \textbf{Usa el almacenamiento en la nube}: Sube tus fotos y vídeos a la nube (iCloud, Google Photos, etc.) y bórralos del teléfono para liberar espacio.
  \item \textbf{Usa una tarjeta microSD si el teléfono lo soporta}: Si tu teléfono tiene ranura microSD, inserta una tarjeta para ampliar el almacenamiento. Úsala para fotos, vídeos y archivos grandes.
  \item \textbf{Borra los archivos descargados que no necesites}: Ajustes > Almacenamiento > Archivos. Revisa y borra los archivos descargados que ya no necesites (documentos, ZIP, APK).
  \item \textbf{Borra el historial del navegador}: El historial del navegador y la caché ocupan espacio. Bórralos periódicamente en Ajustes del navegador.
\end{itemize}

\booksubsection{Usa “Optimizar almacenamiento” para la música y los vídeos}

\begin{itemize}
  \item \textbf{Google Family Link}: En la app de Música o Vídeo, activa "Optimizar almacenamiento" para que el teléfono borre automáticamente los contenidos que no hayas reproducido en mucho tiempo.
  \item \textbf{Revisa el almacenamiento una vez al mes}: Ajustes > Almacenamiento. Revisa qué ocupa más espacio y actúa en consecuencia. Un teléfono con espacio libre funciona más rápido.
  \item \textbf{XI. Diez habilidades de comunicación}El teléfono es, ante todo, una herramienta de comunicación. Estos diez consejos mejoran tu comunicación:
  \item \textbf{Usa "mensajes de voz" para explicaciones largas}: Si una explicación es demasiado larga para escribirla, envía un mensaje de voz en WeChat. Pero no excedas los 60 segundos; si es más largo, llama.
\end{itemize}

\begin{center}
\begin{tabularx}{\linewidth}{|X|X|X|}
\hline
\textbf{Usa "videollamadas" para comunicaciones importantes} & \textbf{: Para comunicaciones importantes (hablar con los padres, reuniones de trabajo), usa videollamadas en lugar de llamadas de voz. La comunicación cara a cara es más efectiva.} & \textbf{Aprende a usar "compartir pantalla" en las videollamadas} \tabularnewline
\hline
Usa el almacenamiento en la nube & : En WeChat, Zoom y otras apps de videollamada, puedes compartir tu pantalla para mostrar a la otra persona lo que ves en tu teléfono. Útil para soporte técnico remoto. & Usa "grupos" de WeChat para comunicaciones de equipo \tabularnewline
Usa una tarjeta microSD si el teléfono lo soporta & : Crea grupos de WeChat para tu equipo o familia. Usa "notas de grupo" para anuncios importantes y "archivos de grupo" para compartir documentos. & Aprende a usar " WeChat Mini Programs" \tabularnewline
: Los Mini Programs de WeChat son apps ligeras que se abren dentro de WeChat. Úsalos para servicios como pedir comida, pagar facturas o reservar citas sin instalar apps adicionales. & Usa "pagos móviles" para todo & : Alipay y WeChat Pay te permiten pagar casi todo sin efectivo. Configura el pago biométrico para mayor comodidad y seguridad. \tabularnewline
Usa "escaneo de código QR" para todo & Borra el historial del navegador & : El escaneo de código QR te permite añadir amigos, pagar, conectarse al Wi-Fi, abrir enlaces y mucho más. Aprende a usar el escáner integrado del teléfono. \tabularnewline
Aprende a usar "compartir ubicación" & : En WeChat, puedes compartir tu ubicación en tiempo real con tus contactos. Útil para quedar con amigos o para que tu familia sepa dónde estás. & Usa " notas compartidas" para colaborar \tabularnewline
\hline
\end{tabularx}
\end{center}

\begin{tcbwarn}{: Las notas compartidas de iCloud o las notas colaborativas de Google Keep permiten que varias personas editen el mismo documento. Útil para listas de la compra, planificación de viajes y notas de reuniones.}
Aprende a usar “traducción automática”

\end{tcbwarn}

\booksection{: La app de Traducción del iPhone y Google Translate permiten traducir texto, voz y conversaciones en tiempo real. Útil para viajar al extranjero o comunicarte con personas que hablan otros idiomas.}

XII. Diez habilidades de aprendizaje\textbf{El teléfono es también una poderosa herramienta de aprendizaje. Estos diez consejos te ayudan a aprender con el teléfono:}:

\begin{itemize}
  \item \textbf{Usa apps de aprendizaje de idiomas}: Apps como Duolingo, Babbel y Memrise te permiten aprender idiomas en cualquier momento. 15 minutos al día son suficientes para progresar.
  \item \textbf{Escucha podcasts para aprender sobre la marcha}: Los podcasts son una gran forma de aprender mientras conduces, caminas o haces tareas. Apple Podcasts, Xiaoyuzhou y otras apps ofrecen miles de podcasts gratuitos.
  \item \textbf{Usa apps de lectura para leer libros}: Apps como WeChat Reading, Kindle y Apple Books te permiten leer libros en el teléfono. Ajusta el tamaño del texto y el brillo para no cansar la vista.
\end{itemize}

\textbf{Mira vídeos educativos en Bilibili/YouTube}\par

: Bilibili y YouTube tienen miles de vídeos educativos gratuitos: tutoriales, cursos, documentales. Úsalos para aprender nuevas habilidades.

\booksection{Usa apps de toma de notas}

: Apps como Notion, Evernote y Apple Notes te permiten tomar notas en cualquier momento. Organiza tus notas con etiquetas y carpetas para encontrarlas fácilmente.

\begin{itemize}
  \item \textbf{Protección de privacidad}Usa apps de flashcards para memorizar
  \item \textbf{: Apps como Anki y Quizlet usan la repetición espaciada para ayudarte a memorizar vocabulario, datos y conceptos de forma eficiente.}Usa apps de gestión de tareas
  \item \textbf{: Apps como Todoist, TickTick y Apple Reminders te ayudan a gestionar tus tareas. Crea listas de tareas, establece plazos y recordatorios.}Usa apps de calendario
  \item \textbf{: Sincroniza tu calendario con el teléfono (Google Calendar, Apple Calendar). Establece alertas para no olvidar citas y reuniones.}Usa apps de meditación
\end{itemize}

\begin{tcbpractice}{: Apps como Headspace, Calm y Xiao Sleep te guían en meditaciones para reducir el estrés y mejorar el sueño.}
XIII. Diez habilidades de vida

\begin{itemize}
  \item El teléfono puede simplificar muchas tareas cotidianas. Estos diez consejos te ayudan a mejorar tu vida:
  \item Usa apps de mapas para navegar
  \item : Gaode Maps y Baidu Maps te permiten navegar en coche, a pie y en transporte público. Descarga los mapas offline para usarlos sin conexión.
  \item Usa apps de transporte
  \item : Apps como Gaode y Baidu Maps muestran horarios de autobús, metro y tren. Úsalas para planificar tus trayectos.
\end{itemize}

\end{tcbpractice}

\booksection{Usa apps de comida a domicilio}

: Meituan, Ele.me y otras apps te permiten pedir comida a domicilio. Compara precios y lee reseñas antes de pedir.

\begin{itemize}
  \item \textbf{Usa apps de compra online}: Taobao, JD.com, Pinduoduo y otras apps te permiten comprar casi todo. Compara precios y lee reseñas antes de comprar.
  \item \textbf{Usa apps de banking online}: La app de tu banco te permite consultar saldos, transferir dinero y pagar facturas. Activa la autenticación de dos factores para mayor seguridad.
  \item \textbf{Usa apps de salud}: La app de Salud del iPhone y Google Fit registran tus pasos, ritmo cardíaco y sueño. Úsalas para controlar tu actividad física.
  \item \textbf{Usa apps de meteorología}: Mira el tiempo por la mañana para decidir la ropa y el paraguas. Activa las alertas de lluvia para no pillarte fuera.
\end{itemize}

\begin{tcbinfo}{Usa apps de reservas}
: Apps como Meituan y Dianping te permiten reservar restaurantes, hoteles y entradas de cine. Úsalas para evitar colas y decepciones.

\end{tcbinfo}

\booksection{Usa apps de traducción}

: La app de Traducción del iPhone y Google Translate permiten traducir texto, voz y conversaciones en tiempo real. Útil para viajar al extranjero.

\begin{itemize}
  \item \textbf{Usa apps de scanner de documentos}: El iPhone tiene un escáner de documentos integrado en Notas. Android también tiene apps de escaneo (como CamScanner); digitaliza documentos en cualquier momento.
  \item \textbf{XIV. Diez habilidades de entretenimiento}El teléfono es también una gran herramienta de entretenimiento. Estos diez consejos te ayudan a disfrutar más:
  \item \textbf{Usa apps de streaming de música}: NetEase Cloud Music, QQ Music y Spotify te permiten escuchar millones de canciones. Crea playlists y descarga música para escuchar sin conexión.
\end{itemize}

\begin{tcbwarn}{Usa apps de streaming de vídeo}
: iQiyi, Tencent Video, Youku y YouTube te permiten ver películas, series y programas. Descarga vídeos para verlos sin conexión.

\end{tcbwarn}

\begin{tcbsummary}{Resumen del capítulo}
\begin{itemize}
  \item Usa apps de podcasts
  \item : Apple Podcasts, Xiaoyuzhou y otras apps ofrecen miles de podcasts gratuitos. Suscríbete a tus favoritos y descarga episodios para escuchar sin conexión.
  \item Usa apps de lectura
  \item : WeChat Reading, Kindle y Apple Books te permiten leer libros en el teléfono. Ajusta el tamaño del texto y el brillo para no cansar la vista.
  \item Usa apps de juegos
  \item : App Store y Google Play tienen miles de juegos. Juega con moderación; establece un límite de tiempo para no excederte.
  \item Usa apps de fotos
\end{itemize}

\end{tcbsummary}
